% 작업장 공통 서식 (W2 이후). tectonic으로 컴파일.
\documentclass[11pt,a4paper]{article}
\usepackage{fontspec}
\usepackage{xeCJK}
\setCJKmainfont{AppleMyungjo}
\setCJKsansfont{Apple SD Gothic Neo}
\setCJKmonofont{Apple SD Gothic Neo}
\xeCJKsetup{CJKspace=true}
\usepackage{amsmath,amssymb,amsthm,mathtools}
\usepackage[margin=23mm]{geometry}
\usepackage{booktabs,tabularx,array}
\usepackage{enumitem}
\usepackage{xcolor}
\usepackage[most]{tcolorbox}
\usepackage[hidelinks]{hyperref}
\linespread{1.25}
\setlength{\parindent}{0pt}
\setlength{\parskip}{0.55em}

\definecolor{ink}{HTML}{1B2433}
\definecolor{acc}{HTML}{2C55C9}
\definecolor{hot}{HTML}{B8650A}
\newtcolorbox{thmbox}{colback=acc!4,colframe=acc!60,boxrule=0.5pt,arc=1pt,left=6pt,right=6pt,breakable}
\newtcolorbox{bridge}{colback=hot!5,colframe=hot!70,boxrule=0.5pt,arc=1pt,left=6pt,right=6pt,breakable}
\newtcolorbox{openprob}{colback=ink!3,colframe=ink!45,boxrule=0.5pt,arc=1pt,left=6pt,right=6pt,breakable}

\newcommand{\blank}{\rule[-0.3em]{0pt}{1.6em}}
\newcolumntype{Y}{>{\raggedright\arraybackslash}X}
\newcommand{\steptable}[1]{\par\noindent%
  \renewcommand{\arraystretch}{1.55}%
  \begin{tabularx}{\linewidth}{@{}p{1.1cm}>{\raggedright\arraybackslash}p{7.6cm}Y@{}}
  \toprule
  단계 & 주장 & 근거 (직접 채울 것)\\ \midrule
  #1
  \bottomrule
  \end{tabularx}}
\newcommand{\Var}{\operatorname{Var}}
\newcommand{\Cov}{\operatorname{Cov}}
\newcommand{\E}{\mathbb{E}}
\newcommand{\Unif}{\operatorname{Unif}}
\newcommand{\round}{\operatorname{rd}}
\newcommand{\sX}{\sigma_X^2}
\newcommand{\sN}{\sigma_N^2}

\title{\sffamily 작업장 W5\\[2pt]\large 조립: 사다리 C단계의 정리와 남은 열린 문제}
\author{}
\date{2026-11-22}
\begin{document}
\maketitle
\vspace{-2.5em}
\begin{center}\small
선행: W1(역정리 사슬), W2(달성), W3(기준선), W4(찌꺼기 항). 사다리: \texttt{draft.md} 정리 골격 v1, \texttt{SPEC.md} P절.
\end{center}

\textbf{사용법.} 이 작업장은 앞의 네 작업장을 하나의 정리로 묶고, 그 정리가 \emph{어디까지} 증명됐고 \emph{어디서부터} 연구인지를 정확히 표시한다. 1–3부는 채우는 표, 4부는 열린 문제다.

\section*{0. 지금까지 가진 것}
\begin{center}\small
\renewcommand{\arraystretch}{1.3}
\begin{tabularx}{\linewidth}{@{}lYl@{}}
\toprule
출처 & 내용 & 성립 범위\\ \midrule
W1 2부 & $\log|G|\ge H(Z\mid W)=I(X;Z\mid W)+H(Z\mid X,W)$, \ $I(X;Z\mid W)\ge\frac12\log(\sX/D)$ & 모든 부호\\
W3 4부 & 불편 복호기면 $I(X;Z\mid W)\ge\frac12\log(1+\sX/D)$ & 모든 불편 부호\\
W3 & 제약 없는 기준선: 하행 $\frac12\log(1+\sX/D)$, 작업자당 상행 $r=\frac12\log\frac{KD}{KD-\sN}$ (불편) & 모든 부호\\
W4 (a)(b)(c) & $H(S\mid X,W)\gtrsim\frac12\log K+\log\frac{\sigma_N}\Delta$, \ $H(S\mid W)\gtrsim\log K+\log\frac{\sigma_X}\Delta$ & 부호기군 $\mathcal F_\Delta$\\
W2 & $\log M\le\log K+\log\frac{\sigma_X}\Delta+\log(2a)+o(1)$, \ $D\le\frac{\sN+\Delta^2/12}K+\varepsilon(a)$ & 방식 $\mathcal S(\Delta,M)$\\
\bottomrule
\end{tabularx}
\end{center}

\section*{1부. 정리 C (사다리 C단계)}
\begin{thmbox}
\textbf{정리 C.} 부호기군 $\mathcal F_\Delta$(작업자별 독립 디더, 눈금 $\Delta$)와 임의의 $M$, 임의의 복호기 $g(Z,W)$에 대해
\begin{align*}
\text{(i)}\quad R_{\mathrm{down}}&\ge\log\frac{K+1}2+\frac12\log\frac{2\pi e(\sX+\sN/K)}{\Delta^2}-H(J),\\
\text{(ii)}\quad R_{\mathrm{down}}&\ge\frac12\log\frac{\sX}D+\frac12\log K-1+\frac12\log\frac{2\pi e\sN}{\Delta^2}-H(J).
\end{align*}
불편 복호기면 (ii)의 첫 항을 $\frac12\log(1+\sX/D)$로 바꿀 수 있다. 또 W2에 의해 같은 $\Delta$에서 $R_{\mathrm{down}}\le\log K+\log(\sigma_X/\Delta)+O(\log\log K)$가 달성된다.
\end{thmbox}

\steptable{
1.1 & $R_{\mathrm{down}}=\log M\ge H(Z\mid W)$ & \\
1.2 & (i): $H(Z\mid W)\ge H(S\mid W)-H(J)$와 W4(b) & \\
1.3 & (ii): $H(Z\mid W)=I(X;Z\mid W)+H(Z\mid X,W)$ & \\
1.4 & $I(X;Z\mid W)\ge\frac12\log(\sX/D)$ & \\
1.5 & $H(Z\mid X,W)\ge H(S\mid X,W)-H(J)$와 W4(a) & \\
}

\section*{2부. 해상도를 목표에 맞추면: $\tfrac12\log K$ 간극}
\textbf{따름정리 C$'$.} 목표 $D\approx\theta\sN/K$에 W2가 쓰는 눈금 $\Delta=\sqrt{12(\theta-1)}\,\sigma_N$을 두면, 부호기군 $\mathcal F_\Delta$ 안에서
\[
R_{\mathrm{down}}=\log_2K+\tfrac12\log_2\gamma+O(\log\log K),
\qquad
R_{\mathrm{down}}-\underbrace{\tfrac12\log_2\!\big(1+K\gamma/\theta\big)}_{\text{제약 없는 하행}}=\tfrac12\log_2K+O(\log\log K).
\]

\steptable{
2.1 & (i)에서 $\log\frac{K+1}2\ge\log K-1$, \ $\frac12\log\frac{2\pi e\sX}{\Delta^2}=\frac12\log\frac{2\pi e\gamma}{12(\theta-1)}$ & \\
2.2 & 따라서 하한 $=\log K+\frac12\log\gamma-c_1(\theta)-H(J)$ & (대수)\\
2.3 & W2 따름정리에서 상한 $=\log K+\frac12\log\gamma+c_2(\theta)+O(\log\log K)$ & \\
2.4 & 제약 없는 하행 $\frac12\log(1+K\gamma/\theta)=\frac12\log K+\frac12\log\gamma-\frac12\log\theta+o(1)$ & \\
2.5 & 빼면 간극 $\frac12\log K+O(\log\log K)$ & (대수)\\
}

\textbf{연습 2.6} $c_1(\theta)$, $c_2(\theta)$를 구체적으로 쓰라. $\theta=2$에서 숫자로. \vspace{3em}

\begin{center}\small
\renewcommand{\arraystretch}{1.3}
\begin{tabular}{rccccc}
\toprule
$K$ & 하한 (i) & 하한 (ii, 불편) & 상한 W2 ($a=12$) & 상한 (가우시안 근사) & 제약 없는 하행\\ \midrule
4 & 1.74 & 1.37 & 5.67 & 3.88 & 0.79\\
16 & 3.39 & 2.93 & 6.87 & 5.40 & 1.58\\
64 & 5.29 & 4.80 & 8.81 & 7.23 & 2.52\\
256 & 7.26 & 6.77 & 10.80 & 9.18 & 3.51\\
1024 & 9.26 & 8.76 & 12.79 & 11.17 & 4.50\\
\bottomrule
\end{tabular}\\[2pt]
$\sX=\sN=1$, $\theta=2$, $\Delta=\sqrt{12}$, $D=2/K$, 비트. $H(J)$는 $a=12$에서 $10^{-5}$ 이하라 뺐다. 하한·상한은 $K$가 4배일 때 약 2비트, 제약 없는 하행은 1비트 오른다.
\end{center}

\textbf{연습 2.7} 표에서 (i)이 (ii)보다 큰 이유를 W4 연습 2.8의 답과 연결해 설명하라. \vspace{3em}

\section*{3부. 논문에 쓸 수 있는 진술}
\begin{bridge}
\textbf{지금 증명된 것 (C단계).} “작업자별 독립 디더를 쓰는 균일 양자화 군 합 부호는, 목표 왜곡 $\theta\sN/K$에 맞춘 해상도에서 하행 레지스터에 $\log_2K+\frac12\log_2\gamma+O(\log\log K)$비트가 필요충분하다. 이는 제약 없는 CEO보다 $\frac12\log_2K$비트 많다.”

\textbf{이것만으로 부족한 이유.} 이 진술은 (1) 부호기군을 $\mathcal F_\Delta$로 한정하고, (2) 해상도 $\Delta$를 목표에 맞춘다고 가정한다. 심사자는 “다른 부호기나 다른 $\Delta$로 피할 수 있지 않은가”를 물을 것이다. 4부가 그 질문들이다.
\end{bridge}

\section*{4부. 열린 문제 (여기서부터 연구)}
\begin{openprob}
\textbf{O1. 해상도–왜곡 연결.} 부호기군 $\mathcal F_\Delta$에서 임의의 복호기가 $D\le\theta\sN/K$를 내려면 $\Delta=O(\sigma_N)$이어야 함을 보여라. 즉 $\mathrm{mmse}(X\mid Z,W)\ge D_{\min}+c\,\Delta^2/K$.\\
\emph{시도와 실패:} W4처럼 $N'$을 조건으로 두면 칸 길이가 $\Delta/K$가 되어 $c\Delta^2/K^2$밖에 안 나온다. $N'$을 모르는 상태의 평균을 다뤄야 한다. $K=1$부터 해 볼 것.

\textbf{O2. 일반 스칼라 부호기 (B단계 = 보조정리 L).} $f_k(Y_k,W)\in\mathbb Z_M$가 임의일 때 $H(Z\mid X,W)\ge\frac12\log K-c$를 $D$와 연결해 보여라.\\
\emph{살아남는 도구:} 보조정리 B(공통 성분 $\bar G$의 분리)는 부호기와 무관하다. 무너지는 것: 보조정리 C의 단조성. 후보: $t\mapsto f_k(c_k+t,W)$의 “$t$에 대한 민감도”를 왜곡으로 하한하고, 군 $\mathbb Z_M$ 위 독립 합의 엔트로피 성장(이산 EPI, 합집합 엔트로피; \texttt{papers/modern/Tao\_sumset\_entropy.pdf} 등)으로 넘긴다. 부분군에 갇히는 경우(D1)를 따로 처리해야 한다.

\textbf{O3. 블록 부호 (A단계).} 찌꺼기가 $\frac n2\log K$로 자라는가. 보조정리 D의 격자판(디더가 보로노이 칸 위 균일, 경계가 $\mathbb R^n/\Lambda$ 위에 흩어짐)이 후보.

\textbf{O4. 공유 디더의 교환 (새 현상 후보, D3).} W4 연습 4.1: 디더를 공유하면 찌꺼기는 줄지만 양자화 오차가 평균되지 않는다. “찌꺼기 $\times$ 왜곡”의 교환 곡선을 정확히 구하면, 저율 영역에서 대가가 줄어드는 조건이 나올 수 있다.
\end{openprob}

\textbf{연습 4.1 (연구 첫걸음).} O1을 $K=1$에서 풀어 보라: $\mathrm{mmse}(X\mid q_1,U_1)\ge D_{\min}+c\Delta^2$인가? ($Y_1=X+N_1$, 빼는 디더 양자화. 힌트: $(q_1,U_1)$이 주어지면 $Y_1$은 길이 $\Delta$인 구간 안에 있다. 그 구간이 $\sigma_N$보다 넓으면?)
\vspace{6em}

\newpage
\section*{부록 A. 힌트}
\begin{itemize}[leftmargin=2.2em]
\item[1.1] $Z$는 $M$개 값 중 하나. 조건부 엔트로피 $\le$ 엔트로피 $\le\log M$.
\item[1.4] W1 2부의 (d)–(j)를 $M$ 대신 $(Z,W)$로. $\hat X=g(Z,W)$.
\item[2.3] W2 따름정리 W2$'$.
\item[2.4] $K\to\infty$에서 $\log(1+x)=\log x+o(1)$.
\end{itemize}

\section*{부록 B. 답}
\begin{itemize}[leftmargin=2.2em]
\item[2.6] $c_1(\theta)=1-\frac12\log_2(2\pi e)+\frac12\log_2(12(\theta-1))$; $\theta=2$면 $1-2.047+1.792\approx0.745$. \ $c_2(\theta)=-\frac12\log_2(12(\theta-1))+\log_2(2a)$; 고정 $a=12$, $\theta=2$면 $-1.792+4.585\approx2.79$. 하한·상한의 상수 차이는 약 3.5비트이고, 대부분은 W2의 꼬리 여유($\log_2 2a$)에서 온다.
\item[2.7] (i)은 $H(Z\mid W)$ 전체를 바로 하한한다(신호 $X$의 폭 $\sigma_X$를 칸 $2\Delta/(K+1)$로 재므로 $\log K$ 전체). (ii)는 $I(X;Z\mid W)$를 왜곡만으로 하한하는데($\frac12\log(1/D)$, 이미 $\frac12\log K$ 포함), 그 과정에서 상수를 조금 잃는다. 차수는 같다.
\item[4.1] (방향만) $\Delta\gg\sigma_N$이면 구간 정보가 $Y_1$을 거의 고정하지 못해 $X$의 추정 오차에 $\Delta^2/12$ 수준이 남는다는 것을 보이면 된다. 엄밀한 하한은 열린 문제로 남긴다.
\end{itemize}
\vspace{1em}
\small\textcolor{ink!60}{표의 수치는 Python으로 계산함 (scratchpad/w25nums.py, w4check.py).}
\end{document}
